% © 2026 Ing. David Jaroš — CC BY-NC-ND 4.0
%
% This work is licensed under a Creative Commons Attribution-NonCommercial-NoDerivatives
% 4.0 International License (CC BY-NC-ND 4.0).
%
% License History: Earlier drafts (up to v0.3) were released under CC BY 4.0.
% From v0.4 onward, all material is released under CC BY-NC-ND 4.0 to protect
% the integrity of the theoretical work during ongoing academic development.

% UBT-AI-PROVENANCE-BEGIN
% schema: ubt-ai-provenance/v1
% tier: C_working
% ai_assistance: disclosed
% human_review: risk-based
% editorial_responsibility: Ing. David Jaroš
% policy: ../../AI_PROVENANCE.md
% notice: Working material; exhaustive human review is not claimed.
% UBT-AI-PROVENANCE-END

\documentclass[11pt]{article}
\usepackage{amsmath,amssymb,amsthm,geometry,hyperref}
\geometry{margin=1in}

\newtheorem{definition}{Definition}
\newtheorem{proposition}{Proposition}
\newtheorem{remark}{Remark}

\title{Gray-Code Transport Layer in UBT Layer 2:\\
Formal Definitions, One-Hot Analysis, and\\
CMB Path-Fingerprint Test}
\author{Ing.\ David Jaro\v{s}}
\date{2026-04-27}

\begin{document}
\maketitle

\begin{abstract}
We formalize the Gray-code transport hypothesis (Layer~2 Transport, L2T) as a
research-track complement to the canonical Hamming $(8,4,4)$ storage fingerprint
(Layer~2 State, L2S). Hamming coding protects state identity; Gray coding
minimizes sequential transition cost. The two schemes are complementary and
address orthogonal questions. We define both layers precisely, analyze
one-hot triqubit color states under each scheme, distinguish the roles of
off-diagonal and diagonal SU(3) generators, and propose a falsifiable
CMB path-fingerprint test based on the \emph{gray\_adjacency\_score}.
\end{abstract}

\tableofcontents
\newpage

% -------------------------------------------------------------------
\section{Layer 2 Structure: State and Transport Sub-layers}

UBT Layer~2 (Coding and Stability) stabilizes admissible field configurations
through discrete selection rules. This layer is subdivided:

\begin{itemize}
\item \textbf{L2S (State Layer):} discrete stabilizer selection rules protecting
the \emph{identity} of admissible configurations.
Canonical fingerprint: extended Hamming $(8,4,4)$ code.
Observable: $P_0$ (syndrome-zero fraction).

\item \textbf{L2T (Transport Layer):} transition-cost rules governing sequential
phase-symbol evolution along the $\psi$-time direction.
Research-track hypothesis: Gray-code ordering.
Observable: $A_{\mathrm{gray}}$ (Gray adjacency score).
\end{itemize}

\begin{remark}[Complementarity]
L2S and L2T address orthogonal questions.
L2S asks: \emph{does this configuration survive?}
L2T asks: \emph{is this path consistent?}
They probe distinct statistical properties of the same observable sequence
and are independently falsifiable.
Gray code does \textbf{not} replace Hamming.
\end{remark}

% -------------------------------------------------------------------
\section{Definitions}

\subsection{Hamming Distance and the Hamming $(8,4,4)$ Code}

\begin{definition}[Hamming distance]
For two binary strings $\mathbf{u}, \mathbf{v} \in \{0,1\}^n$, the
\emph{Hamming distance} is
\[
  d_H(\mathbf{u},\mathbf{v}) = \bigl|\{i : u_i \neq v_i\}\bigr|.
\]
\end{definition}

\begin{definition}[Extended Hamming $(8,4,4)$ code]
The extended Hamming code $\mathcal{H}(8,4,4)$ is a linear binary code with
parameters:
\begin{itemize}
\item block length $n = 8$,
\item dimension $k = 4$ (four information bits),
\item minimum distance $d_{\min} = 4$.
\end{itemize}
Its parity-check matrix $H \in \{0,1\}^{4\times 8}$ satisfies $H\mathbf{c}^T = \mathbf{0}$
for every codeword $\mathbf{c} \in \mathcal{H}$.
The code detects all errors of weight $\leq 3$ and corrects all errors of weight
$\leq 1$.
\end{definition}

\begin{definition}[Hamming syndrome and $P_0$]
Given an 8-symbol phase block $\mathbf{s}_k = (s_{k,1},\dots,s_{k,8}) \in
\{0,1\}^8$, the \emph{syndrome} is
\[
  \mathbf{r}_k = H \mathbf{s}_k^T \in \{0,1\}^4.
\]
The \emph{fingerprint statistic} is
\[
  P_0 = \frac{1}{K}\sum_{k=1}^{K} \delta(\mathbf{r}_k, \mathbf{0}),
\]
where $\delta$ is the Kronecker delta. A statistically significant excess of
$P_0$ over a phase-randomized null constitutes a positive detection of L2S.
\end{definition}

\subsection{Gray Adjacency and the Gray Transport Layer}

\begin{definition}[Standard reflected binary Gray code]
The \emph{standard reflected binary Gray code} on $N = 2^b$ symbols is the
sequence $G : \{0,\dots,N{-}1\} \to \{0,1\}^b$ defined recursively by
$G(0) = 0\cdots0$ and
\[
  G(n) = n \oplus \lfloor n/2 \rfloor,
\]
where $\oplus$ denotes bitwise XOR. By construction,
\[
  d_H\bigl(G(n),\, G(n{+}1)\bigr) = 1 \quad \text{for all } n < N{-}1,
\]
and the code is cyclic: $d_H(G(N{-}1), G(0)) = 1$.
\end{definition}

\begin{definition}[Gray adjacency]
Two symbols $s, s' \in \{0,\dots,N{-}1\}$ are \emph{Gray-adjacent} if they are
consecutive in the Gray ordering modulo $N$:
\[
  \bigl|G^{-1}(s) - G^{-1}(s')\bigr| \equiv 1 \pmod{N}.
\]
Equivalently, $d_H(s_{\mathrm{gray}}, s'_{\mathrm{gray}}) = 1$ where
$s_{\mathrm{gray}} = G^{-1}(s)$ and $s'_{\mathrm{gray}} = G^{-1}(s')$.
\end{definition}

\begin{definition}[Gray adjacency score $A_{\mathrm{gray}}$]
Given a sequence of discrete phase symbols $s_1, s_2, \dots, s_m$, the
\emph{Gray adjacency score} is
\[
  A_{\mathrm{gray}} = \frac{1}{m-1}\sum_{i=1}^{m-1}
    \mathbf{1}\bigl[s_{i+1} \text{ is Gray-adjacent to } s_i\bigr].
\]
For a uniform-random i.i.d.\ sequence over $N$ symbols, the expected score
under the null is
\[
  \mathbb{E}[A_{\mathrm{gray}}^{\mathrm{null}}] = \frac{2}{N}
\]
(two Gray-adjacent neighbors out of $N$ symbols for each $s_i$).
\end{definition}

% -------------------------------------------------------------------
\section{One-Hot Triqubit Color States and Hamming Weight-1 Protection}

The one-hot triqubit encoding assigns:
\[
  |r\rangle \mapsto |100\rangle, \quad
  |g\rangle \mapsto |010\rangle, \quad
  |b\rangle \mapsto |001\rangle.
\]
Each color state has \emph{Hamming weight exactly 1} in the 3-qubit register.

\begin{proposition}[Single-bit-flip leakage detection by the weight-1 sector]
\label{prop:hamming_weight1}
Any single bit-flip error on a one-hot state produces a state outside the
one-hot subspace (either weight 0 or weight $\geq 2$), and is therefore
detectable.
\begin{proof}
A one-hot state $\mathbf{e}_j \in \{0,1\}^3$ (weight 1) can suffer:
\begin{itemize}
\item A $0 \to 1$ flip at position $i \neq j$: produces weight-2 state $\mathbf{e}_i +
\mathbf{e}_j \notin \text{one-hot subspace}$.
\item A $1 \to 0$ flip at position $j$: produces the zero vector $\mathbf{0}$
(weight 0, vacuum state) $\notin \text{one-hot subspace}$.
\end{itemize}
In both cases the result is outside the one-hot subspace. \qed
\end{proof}
\end{proposition}

\begin{remark}
This is an occupation-sector property relevant to L2S: the Hamming weight-1
constraint detects every individual computational-basis bit flip $X_i$. It is
not a standard Pauli stabilizer code and does not detect phase-only errors
$Z_i$. The statement is logically separate from the extended Hamming
$(8,4,4)$ test on full 8-symbol blocks.
\end{remark}

% -------------------------------------------------------------------
\section{SU(3) Generator Decomposition: Diagonal vs Off-Diagonal}

SU(3) has 8 generators (Gell-Mann matrices $\lambda_1,\dots,\lambda_8$), which
split naturally under the one-hot triqubit embedding:

\subsection{Off-Diagonal (Color-Changing) Generators}

$\lambda_1, \lambda_2$ (r$\leftrightarrow$g), $\lambda_4, \lambda_5$
(r$\leftrightarrow$b), $\lambda_6, \lambda_7$ (g$\leftrightarrow$b):

\begin{itemize}
\item Each maps one color to a different color: $|r\rangle \leftrightarrow
|g\rangle$, etc.
\item The Hamming distance between a source and target one-hot state is 2
(one $1\to0$, one $0\to1$).
\item These are \emph{state-changing} transitions governed by L2S.
\item They are \emph{not} Gray-adjacent in the phase symbol sense.
\end{itemize}

\subsection{Diagonal (Phase) Generators}

$\lambda_3$ (isospin) and $\lambda_8$ (hypercharge):

\begin{itemize}
\item Each preserves the one-hot color identity while acting as a phase rotation
within a color channel.
\item Sequential phase steps along $\lambda_3$ or $\lambda_8$ are the natural
domain of Gray transport.
\item L2T (Gray adjacency) applies to \emph{phase transitions}, not color
transitions.
\end{itemize}

\begin{remark}[No overstated claims]
Gray transport does not apply uniformly to all 8 SU(3) generators.
The claim is limited: \emph{sequential phase-symbol paths along the diagonal
generator directions prefer Gray-adjacent transitions}.
Off-diagonal transitions are color-changes governed by Hamming L2S.
\end{remark}

% -------------------------------------------------------------------
\section{Observable Distinction}

\begin{center}
\begin{tabular}{lllll}
\hline
Statistic & Layer & Code & Generator sector & What it tests \\
\hline
$P_0$ & L2S & Hamming $(8,4,4)$ & All (8-block parity) & State survival \\
$A_{\mathrm{gray}}$ & L2T & Gray adjacency & Diagonal ($\lambda_3,\lambda_8$) & Path consistency \\
\hline
\end{tabular}
\end{center}

The two statistics are \textbf{independent}: a sequence can exhibit high $P_0$
and low $A_{\mathrm{gray}}$, or vice versa.

% -------------------------------------------------------------------
\section{CMB Path-Fingerprint Test}

\subsection{Input: Discretized Phase Symbols}

Let $a_{\ell m}$ denote complex CMB spherical harmonic coefficients. Define
discretized phase symbols by:
\[
  s_{\ell m} = \left\lfloor \frac{N}{2\pi}\bigl(\arg(a_{\ell m}) + \pi\bigr)
  \right\rfloor \in \{0,\dots,N{-}1\},
\]
with $N = 16$ (4-bit) or $N = 256$ (8-bit). The sequence is formed by
ordering the $(\ell, m)$ pairs along a fixed path (e.g., row-by-row in $\ell$).

\subsection{Test Procedure}

\begin{enumerate}
\item Compute $A_{\mathrm{gray}}^{\mathrm{obs}}$ on the observed symbol sequence.
\item Generate $K \geq 1000$ null sequences by:
  \begin{itemize}
  \item \textbf{Shuffle null}: randomly permute $s_1,\dots,s_m$, preserving the
  marginal distribution but destroying sequential structure.
  \item \textbf{Phase-randomized null}: replace each $\arg(a_{\ell m})$ with an
  independent draw from $\mathrm{Uniform}[0, 2\pi)$, then re-symbolize.
  \end{itemize}
\item Compute $A_{\mathrm{gray}}^{\mathrm{null},j}$ for each null sequence $j$.
\item Report the one-sided $p$-value:
\[
  p = \frac{|\{j : A_{\mathrm{gray}}^{\mathrm{null},j} \geq A_{\mathrm{gray}}^{\mathrm{obs}}\}|}{K}.
\]
\end{enumerate}

\subsection{Interpretation}

\begin{itemize}
\item $p < 0.01$: positive detection of L2T Gray-transport preference.
\item $p \geq 0.05$: no evidence for L2T at this sensitivity.
\item Non-detection of L2T does not affect L2S $P_0$ status.
\item Detection of L2T does not prove QCD or dark matter; it supports the
Gray transport hypothesis in the UBT L2T sector.
\end{itemize}

\subsection{Expected Null Value}

For uniform i.i.d.\ symbols over $N=16$:
\[
  \mathbb{E}[A_{\mathrm{gray}}^{\mathrm{null}}] = \frac{2}{16} = 0.125.
\]
For $N=256$:
\[
  \mathbb{E}[A_{\mathrm{gray}}^{\mathrm{null}}] = \frac{2}{256} \approx 0.0078.
\]
A detection requires $A_{\mathrm{gray}}^{\mathrm{obs}}$ significantly above
these baselines.

% -------------------------------------------------------------------
\section{Conclusion and Status}

The Gray transport layer (L2T) is formalized here as a research-track hypothesis
complementary to the canonical Hamming L2S fingerprint. The two layers address
distinct physical questions (state identity vs.\ path consistency) and are
independently testable.

\begin{itemize}
\item \textbf{L2S status}: \textsc{canonical}.
  Observable $P_0$; code: Hamming $(8,4,4)$.
\item \textbf{L2T status}: \textsc{research track}.
  Observable $A_{\mathrm{gray}}$; code: reflected binary Gray.
\item \textbf{No conflict}: one-hot triqubit occupation-sector logic is preserved.
  Gray applies only to diagonal (phase) transitions; Hamming governs the
  full 8-block state.
\end{itemize}

The Python implementation of the CMB test is in:\\
\texttt{experiments/forensic\_fingerprint/tools/gray\_path\_symbol\_test.py}

\end{document}
